% Abhakseite „Das kann ich“ – Zone nur im Gesamt (\mitzone)
%% TODO Ich-kann-Sätze: Platzhalter wie in den Aufgabendateien
\begin{abhakseite}
\ifdefined\mitzone
\abhakgruppe{Kennst du schon}
\abhak{1}{Potenzen und Brüche rechnen}
\abhak{2}{Ableitungsregeln sicher anwenden, auch Produkt- und Kettenregel}
\abhak{3}{Potenzregel rückwärts}
\abhak{4}{Den Ableitungsgraphen lesen}
\abhak{5}{Steigungsdreieck und Tangente am Graphen}
\abhak{6}{Potenzen und Brüche rechnen – Fehler finden}
\abhak{7}{Potenzen und Brüche rechnen – selbst rechnen}
\fi
\abhakgruppe{Einheit 1 · Stammfunktion}
\abhak{8}{Stammfunktion}
\abhak{9}{Stammfunktion – Fehler finden, Begründen}
\abhakgruppe{Einheit 2 · Hauptsatz}
\abhak{10}{Hauptsatz}
\abhak{11}{Kurvenlänge über eine vorgegebene Integralformel berechnen}
\abhak{12}{Hauptsatz – Fehler finden, Begründen}
\abhakgruppe{Einheit 3 · Der Graphenblick}
\abhak{13}{Der Graphenblick}
\abhak{14}{Der Graphenblick – Fehler finden, Begründen}
\abhakgruppe{Einheit 4 · Integralfunktion}
\abhak{15}{Integralfunktion}
\abhak{16}{Integralfunktionen mit verschiedenen unteren Grenzen als Verschiebung in y-Richtung begründen}
\abhak{17}{Waagerechte Tangente und Wendepunkt einer Integralfunktion an einer doppelten Nullstelle des Integranden begründen}
\abhak{18}{Integralfunktion – Fehler finden, Begründen}
\end{abhakseite}
